\section{Related Work}
\label{sec:related}

Automated SOC alert triage sits at the intersection of machine learning-based
anomaly detection, natural language-driven autonomous agents, and operational
cybersecurity practice. We survey these three dimensions and identify the research
gaps motivating our design.

\subsection{Machine Learning for Alert Prioritization}

Early large-scale automation of SOC workflows deployed tabular statistical
classifiers over structured endpoint and network log features.
Freitas et al.~\cite{freitas2024cgr} introduced Copilot Guided Response (CGR),
deploying Gradient Boosted Decision Trees (GBDT) with Shapley Additive Explanations
(SHAP) across Microsoft enterprise tenants, achieving a macro-F1 score above 0.87
on the GUIDE dataset. The key advantage of GBDT-based approaches lies in their
calibrated probabilistic outputs and computational efficiency ($<$0.04\,s per
incident), making them suitable for high-throughput production environments.
However, CGR and similar statistical models lack semantic comprehension of
kill-chain attack progressions: they classify individual alert feature vectors
without reasoning across correlated multi-step activity sequences.

Garcia et al.~\cite{garcia2025rosca} proposed ROSCA, integrating alert correlation
with MITRE ATT\&CK technique mapping to rank alert severity. ROSCA demonstrated
scalability advantages over graph-based correlation engines but required manually
curated MITRE mappings for each deployment environment.
Freitas et al.~\cite{freitas2026aip} subsequently introduced Adaptive Incident
Prioritization (AIP), dynamically re-ranking queued incidents based on evolving
threat context, reducing analyst queue overflow by approximately 23\% in production.
Ndichu et al.~\cite{ndichu2026survey} provide a comprehensive survey of AI-driven
alert screening approaches, cataloguing 47 methods and identifying that high
false-positive tolerance remains the predominant unresolved challenge across systems.

\subsection{LLM Agents for Autonomous Cyber Investigation}

The emergence of tool-augmented LLMs enabled a qualitative shift from passive
classification to active evidence gathering. Wei et al.~\cite{wei2025cortex}
proposed CORTEX, a multi-agent debate architecture in which independent LLM
instances challenge each other's threat hypotheses before converging on a verdict.
CORTEX improved investigation consensus quality but introduced substantial per-alert
API costs (approximately \$24 per complex multi-hop investigation), limiting
deployment to Tier-2 escalation workflows.

Habibzadeh et al.~\cite{habibzadeh2025llmsoc} survey LLM integration patterns
across SOC workflows, identifying five operational models: prompt-only, RAG-augmented,
tool-using ReAct, multi-agent debate, and hybrid human-AI handoff. Their analysis
finds that tool-using ReAct variants offer the best cost-accuracy tradeoff under
standard enterprise budgets, but are most vulnerable to premature investigation
termination on empty query returns.

Tsegai et al.~\cite{tsegai2026aida} provide the most rigorous empirical characterization
of this vulnerability through AIDA (Automated Intrusion Detection Agent), evaluating five
distinct tool-using agent architectures across diverse enterprise alert
corpora. All five strategies exhibited a true-positive false negative rate of at
least 40.4\%, with the dominant failure mode being premature termination: when an
initial log query returns an empty result set, the LLM incorrectly infers absence
of malicious activity and closes the investigation without further evidence gathering.
This failure mode is structurally absent from pure-classifier approaches and
represents the critical gap exploited in our Evidence Gate design.

\subsection{Benchmark Evaluation Frameworks}

Realistic evaluation of autonomous cyber investigation agents requires multi-hop log
environments reflecting authentic enterprise attack chains. Wu et al.~\cite{wu2026excytin}
introduce ExCyTIn-Bench, comprising 57 relational log tables and 8 complete attack
chains sourced from real-world breach telemetry, with an investigation reward metric
measuring the fraction of attacker footsteps correctly identified. The state-of-the-art
investigation reward plateaus at 0.606, indicating that existing agents correctly trace
only approximately 60\% of attack activities even under ideal laboratory conditions.

Deason et al.~\cite{deason2025cybersoceval} develop CyberSOCEval, extending LLM
evaluation to malware analysis and threat intelligence reasoning tasks. TechniqueRAG,
proposed by Lekssays et al.~\cite{lekssays2025techniquerag}, addresses knowledge
hallucination in ATT\&CK technique annotation by retrieving structured STIX~2.1
objects from curated CTI corpora, demonstrating 19\% improvement in technique
mapping accuracy over zero-shot LLM baselines.

\subsection{Operational Human-in-the-Loop Considerations}

Deployment of automated triage in production SOCs introduces human-factor risks
independent of classification accuracy. Tilbury and Flowerday~\cite{tilbury2026vigilance}
conducted a longitudinal field study across four enterprise SOC deployments, finding
that sustained automation reliance induces vigilance decrements: analysts who routinely
approve automated verdicts without independent verification exhibit a 31\% reduction
in anomaly detection sensitivity over six months. This ``automation complacency''
effect mandates explicit, calibrated human escalation boundaries rather than
unrestricted automation.

Alahmadi et al.~\cite{alahmadi2022falsepositives} established through qualitative
interviews with 20 SOC practitioners that 99\% of raised security alarms are
ultimately classified as non-actionable, with the majority constituting benign
positives (legitimate administrative activity) rather than sensor false positives.
Sadlek et al.~\cite{sadlek2025killchain} demonstrate that severity-based triage
integrating kill-chain attack graph reasoning reduces analyst mean-time-to-respond
(MTTR) by 38\% over threshold-based alert filtering alone.

\subsection{Research Gap and Positioning}

Table~\ref{tab:literature} summarizes our systematic comparison of related literature.
The analysis reveals a consistent tripartite gap: (1) statistical classifiers lack
kill-chain semantic reasoning; (2) LLM agents suffer from the empty-query premature
termination failure mode; and (3) no existing system jointly addresses probabilistic
calibration, structured knowledge retrieval, bounded investigation planning, and
empirically validated human escalation thresholds. The proposed AI SOC Analyst Agent
addresses all four dimensions through a single integrated spec-driven architecture.

\begin{table}[t]
\caption{Systematic Comparison of Related SOC Triage Literature}
\label{tab:literature}
\centering
\resizebox{\columnwidth}{!}{%
\begin{tabular}{llll}
\toprule
\textbf{System} & \textbf{Year} & \textbf{Core Approach} & \textbf{Key Limitation} \\
\midrule
CGR~\cite{freitas2024cgr}        & 2024 & GBDT + SHAP        & No kill-chain semantics \\
CORTEX~\cite{wei2025cortex}      & 2025 & Multi-LLM Debate   & \$24/alert API cost \\
AIDA~\cite{tsegai2026aida}       & 2026 & Tool ReAct         & $\geq$40.4\% FNR \\
ExCyTIn~\cite{wu2026excytin}     & 2026 & Multi-hop Logs     & Reward plateau 0.606 \\
Vigilance~\cite{tilbury2026vigilance} & 2026 & SOC Field Study  & Complacency trap \\
Alahmadi~\cite{alahmadi2022falsepositives} & 2022 & Field Survey & BP vs.\ FP conflation \\
\textbf{Ours}                    & 2026 & GBDT+RAG+ReAct+Gate & (This work) \\
\bottomrule
\end{tabular}%
}
\end{table}
